\section{Nearly quadratic two-class listing}\label{sec:two-class-fast}
\subsection{A bounded raw moment in the actual simple graph law}
Fix $a\ge3,b\ge1$ or $(a,b)=(2,1)$, put $D=a+b$,
$\delta=a-b$, and choose one fixed odd prime
$p_0\nmid b$. All limits in the next input keep $p_0$ fixed.
For independent degree-$a$ internal configurations and any prescribed
simple $b$-biregular cross matrix $B$,
Appendix~\ref{app:configuration-mixing} and
Corollary~\ref{thm:cubic-raw} give
\begin{equation}\label{eq:raw-config-moment}
 \E_{\rm config}|\ker_{\F_{p_0}}(A-\delta I)|=2p_0+o(1)
\end{equation}
uniformly over $B$, at every admissible $N$.
The appendix treats $a\ge3$ without requiring assortativity. For the
remaining pair
$(2,1)$, the sharp general-frame scalar moment and invariance under
relabeling give exactly the same law for every prescribed matching.
The internal factors remain actual uniform configurations.

Let $K=|\ker_{\F_{p_0}}(A-\delta I)|$. Each configuration contains
the planted $p_0$ vectors. Simplicity of both internal configurations
has probability $s_N\to\exp[-(a^2-1)/2]>0$. Since $K-p_0\ge0$,
\begin{equation}\label{eq:raw-simple-transfer}
 \E_{\rm simple}K
 =p_0+\frac{\E_{\rm config}[(K-p_0)\one_{\mathcal S}]}{s_N}
 \le p_0+\frac{\E_{\rm config}K-p_0}{s_N}.
\end{equation}
In particular
\begin{equation}\label{eq:raw-simple-moment}
 \limsup_N\sup_B\E_{\rm simple}K
 \le\bigl[1+e^{(a^2-1)/2}\bigr]p_0.
\end{equation}
This is an upper bound, not an asserted simple-model limiting value.
Finitely many initial admissible sizes give a finite uniform supremum
over all sizes. Averaging permits any independent random cross law.
No unbounded expectation is transferred through qualitative contiguity.

\begin{theorem}[Full two-class listing in expected nearly quadratic work]\label{thm:fast-full-list}
For every fixed known $a\ge3,b\ge1$, and also for $(a,b)=(2,1)$,
there is an always-correct Las Vegas
algorithm which lists every equitable $(a,b)$ bisection, up to exchange,
on every simple input graph. In the model with independent uniform simple
internal graphs and any prescribed simple cross graph independent
of the internals, its unconditional expected bit work is
$\widetilde O_{a,b}(n^2)$, its expected workspace is
$\widetilde O_{a,b}(n)$ bits, and its expected output size is
$O_{a,b}(n)$ bits, uniformly over the prescribed cross graph.
All admissible block sizes are covered. No real or modular corank-one
promise and no numerical eigengap is used.
\end{theorem}
\begin{proof}
The input domain is simple graphs. Check $D$-regularity; a graph
failing this check has no equitable $(a,b)$ partition, so return the empty list.
Apply Corollary~\ref{fk:cor:moment} to the actual-law bound
\eqref{eq:raw-simple-moment}. On every input, the full base-field
kernel is certified, every coefficient vector is enumerated, and
integer neighbor counts verify every candidate. Every planted input
has a nonempty list. The field extension used by preconditioning
grows with $n$ only for computational success; the probabilistic
kernel moment is at the fixed base prime $p_0$.
\end{proof}

For a random cross factor, Theorem~\ref{thm:boolean} makes this list
unique with probability $1-O_{a,b}(N^{-\nu})$, where
$\nu$ is defined in \eqref{eq:ambiguity-exponent}.
There is no algorithmic error, so this rate has no separate cutoff or
real-corank failure term. Section~\ref{sec:sharp-ambiguity} proves
the exact leading ambiguity and optimal statistical-error constants.
For arbitrary prescribed cross graphs, the real planted-line theorem
proved above gives uniqueness with qualitative probability $1-o(1)$ uniformly
over $B$, without claiming the same quantitative random-cross rate.

\subsection{Unknown assortative degree pair}
If only the total degree $D$ is supplied, the same expected-work
guarantee can be obtained with high-probability recovery of $a,b$.
At $D=3$ the only positive assortative pair is $(2,1)$ and its target
is $\delta=1$. At $D\ge4$, the actual internal degree is at least
three. Enumerate the finitely many positive targets
$\lambda'=D-2b'$ with $1\le b'<D/2$, and choose a fixed odd prime
$p_0>2D$ for every trial. The generic fixed-target configuration
count in Section~\ref{sec:frame} gives
\[
 \E_{\rm config}|\ker_{\F_{p_0}}(A-\lambda'I)|
 =2p_0^{\ell_{p_0}}+o(1),\qquad
 \ell_{p_0}=\nullity_{\F_{p_0}}(Q-\lambda'I).
\]
Here $Q=\left(\begin{smallmatrix}a&b\\b&a\end{smallmatrix}\right)$
has quotient eigenvalues $D,\delta$. The prime choice preserves all
these integer ranks: $\ell_{p_0}=1$ for $\lambda'=\delta$ and zero
for any other candidate. Conditioning gives a bounded raw moment
at every candidate target. Thus the finite collection of full-kernel
listings still costs expected $\widetilde O_D(n^2)$ bit work.
The fixed-target exclusion theorem, by a finite union, leaves only
the planted target with probability $1-o(1)$ uniformly over $B$.
Exact neighbor counts then recover $a,b$. No claim about unproved
cubic moments at other integer shifts is needed.
